% Problema: tres raíces con el teorema de Bolzano (castellano).

\begin{exercise}[Tres raíces con el teorema de Bolzano]
Demuestra que la ecuación $x^3 - 3x + 1 = 0$ tiene exactamente tres
soluciones reales, y localiza cada una en un intervalo de longitud $1$.

\dmarks{3}
\answerspace{6cm}

\begin{hint}
Evalúa el polinomio en los enteros entre $-2$ y $2$ y busca cambios de signo.
\end{hint}

\begin{answer}
Hay una raíz en cada uno de los intervalos $(-2, -1)$, $(0, 1)$ y $(1, 2)$.
\end{answer}

\begin{solution}
El polinomio $p(x) = x^3 - 3x + 1$ es continuo en todo $\mathbb{R}$, y
\[
  p(-2) = -1, \quad p(-1) = 3, \quad p(0) = 1, \quad p(1) = -1, \quad
  p(2) = 3 .
\]
Cambia de signo en $[-2, -1]$, en $[0, 1]$ y en $[1, 2]$, así que el teorema
de Bolzano da una raíz en cada uno de los intervalos abiertos $(-2, -1)$,
$(0, 1)$ y $(1, 2)$. Como son disjuntos, son tres raíces distintas. Y no
puede haber más: un polinomio de grado $3$ tiene a lo sumo tres raíces
reales.
\end{solution}

\begin{marking}
Un punto por evaluar y encontrar los tres cambios de signo, otro por aplicar
Bolzano comprobando la continuidad, y otro por justificar que no hay más de
tres. Es frecuente encontrar solo dos raíces por no evaluar en $-2$.
\end{marking}
\end{exercise}
